\scanpage{20}
\mspage{\handwritten{5}}

distance from both, these three half-circles joining at the two poles $P^+$ and $P^-$. The ``disc'' $Q_i$ around $s_i$ has a conical structure around $s_i$ (as has any conformal pointed disc), and we may take the concentric \deleted{\unsure{disc|dis-|discs}} discs $\handwritten{\lambda}_i Q_i$ with
\[
0\mathrel{\overset{\deleted{\unsure{i}}}{\handwritten{<}}}\handwritten{\lambda}_i\handwritten{<}1.
\]
T\deleted{\unsure{t}}he model I had in mind was that the (non degenerate) admissible systems of discs around the points $s_i$ $(i\handwritten{\in\{}0,1,\infty\handwritten{\}})$ \deleted{\unsure{wh}}are exactly the systems of discs $\handwritten{\lambda}_i Q_i$, with $\handwritten{\lambda}_i$ as above. (If we allow some discs to be degenerate, this means that instead of the unequality above we merely demand $0\handwritten{\leq\lambda}_i\handwritten{<}1$, $0$ not excluded.)

This model, if correct, would give a rather precise description of the inclusion relationships between pants, when these are considered as embedded in the sp\deleted{\unsure{b}}here. The intersection of all would be this system of th\deleted{\unsure{ese}}ree half circles $C_i$, and the two poles $P^+,P^-$ would play a significant role in the geome\handwritten{t}ry of the pants, from this point of view. But it does'nt seem that neither those half circles (which need not be geodesical I guess)\handwritten{,} no\deleted{t}r the two poles\handwritten{,} have ever been described as intrinsically associated to a pant. Of course, this model would give alternative ``parameters'' $\handwritten{\lambda}_i$ for describing a pant, which are best suited for grasping the pants in terms of spherical geometry. The next question would be an understanding of the relationship between these parameters, and Thurston's $\handwritten{\ell}_i$. Maybe it is unreasonable to expect that for given index $i\handwritten{\in\{}0,1,\infty\handwritten{\}}$, the length $\deleted{\handwritten{\lambda}}\handwritten{\ell}_i$ depends only on $\handwritten{\lambda}_i$ and not on the other parameters $\handwritten{\lambda}_j$ - and for this reason, \deleted{\unsure{it is}} the in\handwritten{t}uition at the beginning of this letter, using Thurston's coordinate functions and notably the $\handwritten{\ell}_i$'s to get a fibration structure on $\widetilde{\mathrm{MB}}$ over $\widetilde{\mathrm M}$, in terms of a given decomposition of $X_{\unsure{o|0}}$ into pants, is probably not really relevant, namely it is non intrinsic. Assuming the model I am suggesting is correct, the accurate description \deleted{\unsure{but}} of $\widetilde{\mathrm{MB}}$ in terms of $\widetilde{\mathrm M}$ would be
\[
\widetilde{\mathrm{MB}}\handwritten{\simeq}\widetilde{\mathrm M}\mathbin{\mathrm{x}}\handwritten{[}0,1\handwritten{[}^{I},
\]
where the second factor on the right hand sight refers to the system of multipliers $\handwritten{\lambda}_i$ $(i\handwritten{\in} I)$, tied to the $r_i$ above by $r_i=\handwritten{\lambda}_i\handwritten{\rho}_i$.

My question of course is whether you have any information or idea to propose, especially on the basic problem of relying pants to spherical geometry, and more specifically, \deleted{\unsure{if}}\handwritten{wh}ether the model above is likely to hold, or is defintely false. Also, one difficulty we found with hyperbolic geometry of conformal surfaces, \inserted{\handwritten{is}} that apart from existence and unicity of the hyperbolic structure (compatible with the given conformal one and for which the boundary is geodesic), there seems to be little
% archive sheet number: 19
